\section{未知内容防线：Predictive Detection、校准与人工复核}

生成式工具与新的上传内容都会产生 reference set 未见过的对象，因此下一层是 predictive detection。这里的关键不是“模型能不能给分”，而是怎样把 score 变成可操作且可承受的审核队列：阈值要随风险与容量变化，输出要经过 calibration，最终决策必须保留 human review 和申诉边界。

\subsection{从 classifier score 到 verified outcome}

\term{predictive classifier}（预测分类器）根据训练数据学习输入与风险标签之间的统计关系，输出类别或 score。\term{human review}（人工复核）是由经过培训、受权限与健康保护约束的人员根据上下文和 policy 做出决策，而不是简单“看一眼模型结果”。模型输出是 triage signal，不是法律结论。

\lecturefigure{04-predictive-classifier.png}{Predictive AI 把覆盖范围扩展到已知 hash 之外，但必须经过队列和复核。}{本地字幕 10:10--13:40；Thorn Safer Predict 官方机制；概念重绘。}

\begin{knowledgebox}{读图：Score 的含义取决于训练和校准}
模型先产生 raw score，平台再依据 threshold 将对象分入不同优先级；reviewer 看到的是经过权限过滤的 case context，并将 verified outcome 写回受控系统。Raw score 的大小不能跨 model version、媒体类型或客户环境直接比较。部署前应明确 score 是 ranking、estimated probability 还是 uncalibrated margin，并保存 model/version 与 threshold/version。
\end{knowledgebox}

\teachervoice{Cordua 用 Safer Match 到 Safer Predict 的演进说明：当攻击者或生成模型不断制造新对象时，只查已知 hash 已经不够；但“Predict”这个名字也不意味着系统可以跳过人。课堂的真实重点是扩大候选覆盖并把最值得看的项提前。}

\subsection{Precision、recall、calibration 与 queue capacity}

本节把模型指标和运营约束放在同一张图里。\term{precision}（精确率）表示被系统标为正例的对象中有多少经验证为正例；\term{recall}（召回率）表示所有真实正例中有多少被系统找到；\term{calibration}（校准）表示预测概率与长期实际发生频率是否一致。三者都不能脱离 base rate、label process 与 reviewer capacity 单独优化。

\lecturefigure{10-evaluation-queue.png}{模型质量与复核能力必须一起设计，阈值是运营决策而非纯模型常数。}{本地字幕 10:10--16:20；NIST AI RMF 风险测量思想；概念重绘。}

\begin{importantbox}{指标定义与符号解释}
记真正例、假正例、假负例分别为 $TP,FP,FN$，则
\[
\mathrm{precision}=\frac{TP}{TP+FP},\qquad
\mathrm{recall}=\frac{TP}{TP+FN}.
\]
其中 precision 决定审核员在告警中遇到多少无效项，recall 决定系统遗漏多少真实风险。若每天候选数为 $N$、阈值下告警率为 $q(\tau)$、每名审核员日容量为 $c$、人数为 $R$，稳定运行至少要求
\[
Nq(\tau) \lesssim Rc,
\]
其中 $\tau$ 是阈值。这个不等式是容量约束，不代表安全目标只需“让队列不积压”。
\end{importantbox}

\begin{warningbox}{高 accuracy 可能掩盖极低 base rate}
在真实正例极少的环境中，一个几乎总预测负例的模型也可能有很高 accuracy。平台应按媒体类型、地区、语言、年龄相关场景和产品入口分层报告 precision/recall，并监控 distribution shift。阈值变化既影响漏报，也会改变审核员暴露量和法定报告负担。
\end{warningbox}

\subsection{Moderator triage 是安全系统的一部分}

人工复核不是无限、免费的最后一步。\term{triage}（分诊/优先级排序）把候选按严重度、置信度、时效性与可行动性排序；好的工具还会隐藏非必要视觉信息、合并重复对象、提供最小上下文并支持轮班与心理健康机制。本节读图时要同时观察“决策质量”和“人类暴露量”，不能只追求吞吐。

\lecturefigure{05-moderator-triage.png}{Triage 应最大化保护价值，同时最小化审核员不必要的敏感暴露。}{本地字幕 13:40--16:20；基于 Safer Review 讨论的概念重绘。}

\begin{knowledgebox}{读图：Case package 优于裸文件队列}
输入层先做 hash、模型排序和重复合并；workbench 只展示执行 policy 所需的最小证据；reviewer decision 与 rationale 被结构化记录；case package 再携带 account、时间、关联事件和 action history。这样既减少重复观看，也让二次复核、质量抽样、报告和 incident review 可以基于同一可审计对象进行。
\end{knowledgebox}

\begin{warningbox}{Reviewer wellbeing 不能靠“心理素质”解决}
组织需要 exposure minimization、培训、轮换、休息、专业支持、访问监控和 clear escalation。模型误报会带来真实的人类成本，因此 false positive 不只是 dashboard 数字；反过来，也不能为了减少暴露而把 threshold 调到系统性漏报。产品、模型和人员政策必须共同承担风险。
\end{warningbox}

\subsection{本章小结}

Predictive detection 的成功标准不是 leaderboard accuracy，而是 calibrated signal 能否在可承受队列中帮助审核员更早处理高风险对象。模型指标、阈值、case packaging、人员健康与申诉机制必须共同验收。

